% Generated draft: requires mathematical and source review.
\documentclass[11pt]{article}
\usepackage[margin=25mm]{geometry}
\usepackage{fontspec}
\setmainfont{Latin Modern Roman}
\usepackage{amsmath,amssymb,mathtools,unicode-math}
\setmathfont{Latin Modern Math}
\usepackage{enumitem,needspace,xurl,hyperref}
\setlength{\parindent}{0pt}
\setlength{\parskip}{5pt}
\setlength{\emergencystretch}{3em}
\newcommand{\N}{\mathbb{N}}
\newcommand{\Z}{\mathbb{Z}}
\newcommand{\Q}{\mathbb{Q}}
\newcommand{\R}{\mathbb{R}}
\newcommand{\C}{\mathbb{C}}
\newcommand{\F}{\mathbb{F}}
\newcommand{\A}{\mathbb{A}}
\newcommand{\PP}{\mathbb{P}}
\newcommand{\E}{\mathbb{E}}
\newcommand{\eps}{\varepsilon}
\newcommand{\dd}{\,\mathrm{d}}
\begin{document}
\section*{UnsolvedMath ID: 30001218}
\textbf{Source problem code:} catalogue entry
\subsection{Definitions and mathematical statement}
Let \(G\) be a complex reductive group acting on a projective variety \(X\), and suppose that the orbit set \(G\backslash X\) is finite. The supplied question asks whether, for every such group, variety, and action, there exists an equivalence of triangulated categories
\[
D^b_{G,c}(X)\simeq\operatorname{Perf}\bigl(\operatorname{Ext}(IC)\bigr).
\]
The source does not specify enough conventions to make this a unique mathematical assertion. In particular, the coefficient system, the precise equivariant constructible derived category denoted by \(D^b_{G,c}(X)\), and the object denoted by \(IC\) must be specified. No smoothness assumption on \(X\) or connectedness assumption on \(G\) is supplied.

One possible interpretation of \(IC\) is a direct sum of intersection cohomology complexes associated with orbit closures. This interpretation still requires specifying which equivariant local systems on the orbits are included, their coefficients, and the normalization of the complexes. The record does not say whether only constant local systems or all relevant irreducible equivariant local systems are intended.

For a specified triangulated category \(\mathcal D\) and a specified object \(IC\in\mathcal D\), its graded self-extension algebra is
\[
E=\bigoplus_{n\in\Z}\operatorname{Hom}_{\mathcal D}(IC,IC[n]),
\]
with multiplication given by shifted composition. A different possible target uses the differential graded algebra
\[
A=\operatorname{RHom}(IC,IC)
\]
in a chosen differential graded enhancement of \(\mathcal D\). For a differential graded algebra \(B\), \(\operatorname{Perf}(B)\) denotes the smallest full triangulated subcategory of its derived category of right differential graded modules that contains the regular module \(B\) and is closed under direct summands. Thus interpreting the target as \(\operatorname{Perf}(E)\), with \(E\) retaining its grading and given zero differential, differs from interpreting it as \(\operatorname{Perf}(A)\). The supplied notation does not resolve this distinction or specify another module convention. Consequently, the displayed equivalence is preserved as the proposed question, with these unresolved qualifications; neither interpretation is asserted to be the uniquely intended formulation.
\subsection{Short English statement}
For every action of a complex reductive group on a projective variety with finitely many orbits, the source asks whether its equivariant constructible derived category is equivalent to a perfect-module category built from intersection cohomology extensions. The coefficients, intersection cohomology object, and graded versus differential graded meaning of the target remain unspecified. The supplied partially solved label is dated curation requiring independent review.
\subsection{Sources}
\begin{itemize}
\item Supplied local registry record: permanent UnsolvedMath ID \(30001218\), source problem code \texttt{catalogue entry}, title ``Derived Equivalence from Intersection Cohomology Extensions,'' in the Oberwolfach Reports Open Problems collection. The proposed year is 2009; no individual proposer is identified.
\item Supplied report citation: \emph{Enveloping Algebras and Geometric Representation Theory} (2009), Oberwolfach Reports, pp.\ 801--866, DOI \texttt{10.4171/owr/2009/15}. \url{https://doi.org/10.4171/owr/2009/15}. The record describes this report as documenting workshop discussions and open problems, but gives no problem-specific page or passage.
\item Supplied literature source, labeled ``Schnurer flag-variety theorem'': \url{https://arxiv.org/abs/0809.4785}.
\item Supplied literature source, labeled ``later equivariant derived-category work'': \url{https://home.mathematik.uni-freiburg.de/soergel/PReprints/EMGR.pdf}.
\item Supplied literature source, labeled ``original OWR report'': \url{https://doi.org/10.4171/owr/2009/15}.
\item Literature assessment dated exactly \texttt{2026-08-21}: the supplied curation labels the problem partially solved and describes the Soergel--Lunts equivalence as proved in substantial cases, including flag varieties, while the broad extracted formulation is described as not established in full generality. This is supplied dated curation, not an independently verified literature assessment; its time-sensitive status and cited claims require independent checking.
\end{itemize}
\end{document}
